% body.tex — P04 paper body
% v4 draft, April 2026

% ===================================================================
\begin{abstract}
Employment trajectory topology, as revealed by persistent homology of
Vietoris--Rips filtrations, captures structural features of career pathways.
Single-parameter persistent homology, however, treats all trajectories
identically regardless of income. We introduce a bifiltration on British
Household Panel Survey and Understanding Society data (1991--2023,
$N = 27{,}280$), where simplices are filtered simultaneously by pairwise
distance in PCA-embedded trajectory space and by an income threshold
parameter. The signed measure decomposition from multi-parameter persistent
homology reveals that $\HH{1}$ topological features --- loops representing
cyclical or multi-pathway career structures --- concentrate at mid-to-high
income levels. Quartiles Q3 and Q4 carry approximately four times the
$\HH{1}$ mass of Q1, and this stratification is significant under 999
income-label permutations ($p < 0.001$). Single-parameter persistent
homology on the same landmarks shows no income gradient in $\HH{1}$ feature
density, confirming that the bifiltration captures structure invisible to
standard topological methods. Complementary invariants --- the Hilbert
function and the rank invariant --- confirm the income asymmetry and reveal
that all $\HH{1}$ features persist to infinity along the income axis, ruling
out income-confined topological structures. We argue that the topological
simplicity of low-income trajectory space --- fewer loops, fewer alternative
pathways --- formalises a structural poverty trap: limited career diversity
constrains escape routes from low-income states.
\end{abstract}

\bigskip\noindent\textbf{Keywords:} multi-parameter persistent homology,
signed measures, bifiltration, poverty traps, topological data analysis,
employment trajectories, BHPS, Understanding Society

% ===================================================================
\section{Introduction}
\label{sec:intro}

\subsection{Motivation: The Limitations of Single-Parameter Topology}

Paper~1 in this programme applied Vietoris--Rips persistent homology to UK
employment trajectories embedded in a 20-dimensional PCA space, finding
non-trivial $\HH{0}$ and $\HH{1}$ structure that exceeds the expectations of
first-order Markov null models. That analysis uses a single filtration
parameter --- the Rips radius --- which treats all pairwise distances equally
regardless of economic outcome. A trajectory cluster in a low-income region
of the space has the same topological weight as one in a high-income region.

Poverty trap theory \citep{azariadis2005poverty, barrett2013economics}
predicts structurally distinct configurations at low income levels ---
absorbing states, threshold effects, bifurcation points. Single-parameter PH
cannot test this prediction because it lacks an income axis. The question
motivating this paper is whether the \emph{topological} structure of
trajectory space varies systematically with income.

\subsection{Multi-Parameter Persistence: From Theory to Computation}

Multi-parameter persistent homology extends single-parameter PH by indexing
the filtration by two or more real-valued parameters simultaneously. The
mathematical theory was established by \citet{carlsson2009theory}, who proved
that no complete barcode decomposition exists in the multi-parameter case ---
unlike single-parameter PH, where the barcode is a complete invariant. This
negative result stalled computational applications for over a decade.

Recent advances have made multi-parameter PH practical. The signed barcode
decomposition \citep{botnan2022signed} and its computational implementation
in the \texttt{multipers} library \citep{loiseaux2023framework, loiseaux2024multipers}
provide tractable invariants --- signed measures, Hilbert functions, rank
invariants --- that capture the essential multi-parameter structure. This
paper is among the first applications of these tools to social science panel
data.

\subsection{Research Questions}

\begin{enumerate}[label=(\arabic*)]
  \item Does the Rips $\times$ income bifiltration reveal topological
    structure absent from single-parameter Rips PH?
  \item How does $\HH{1}$ topological complexity distribute across income
    levels, and is this distribution significantly different from what
    distance structure alone would predict?
  \item What does income-stratified topology tell us about the structure of
    poverty traps?
  \item Do $\HH{1}$ features born at low income persist across the full
    income range, or are they confined to low-income regions of the
    bifiltration?
\end{enumerate}

\subsection{Paper Roadmap}

Section~\ref{sec:lit} reviews the poverty trap and multi-parameter PH
literatures. Section~\ref{sec:methods} describes the bifiltration
construction, signed measure computation, and permutation null testing.
Section~\ref{sec:results} presents results from the 999-permutation
analysis, P01 baseline comparison, GMM regime overlay, and complementary
invariants (Hilbert function and rank invariant).
Section~\ref{sec:discussion} discusses the interpretation of topological
simplicity as a poverty trap signature.

% ===================================================================
\section{Literature Review}
\label{sec:lit}

\subsection{Poverty Traps and Threshold Effects}

Poverty trap theory posits the existence of multiple equilibria in income
dynamics, where self-reinforcing mechanisms hold individuals below critical
thresholds \citep{azariadis2005poverty}. \citet{barrett2013economics}
distinguish structural traps --- where the asset or human capital threshold
is the binding constraint --- from stochastic traps, where shocks push
individuals below a threshold from which recovery is difficult. In the UK
context, \citet{shildrick2012poverty} document persistent cycling between
low-pay and no-pay states that resembles a structural trap without requiring
a sharp threshold.

Markov models of income dynamics \citep{cappellari2004modelling} capture
transition probabilities but not the geometric structure of the trajectory
space. The contribution of topological methods is to characterise the
\emph{shape} of career pathway space at each income level --- not just
transition rates, but the connectivity and cyclical structure of available
pathways.

\subsection{Single-Parameter Persistent Homology for Social Science}

Paper~1 in this programme applied VR PH to UK employment trajectories,
finding significant $\HH{0}$ (cluster) and $\HH{1}$ (loop) structure. The
Markov memory ladder null test confirmed that trajectory topology exceeds
first-order Markov expectations. However, single-parameter PH is
outcome-agnostic: the filtration parameter is pairwise distance in the
embedding, which carries no information about income.

Post-hoc stratification --- computing PH separately on low-income and
high-income subsets --- changes the point cloud and thus the simplicial
complex; it does not decompose the \emph{joint} topology of the full space.
Multi-parameter PH provides a principled framework for this decomposition.

\subsection{Multi-Parameter Persistent Homology}

The theoretical foundations of multi-parameter PH were established by
\citet{carlsson2009theory}, who showed that while persistence modules over a
totally ordered set decompose into interval summands (yielding the barcode),
persistence modules over a partially ordered set --- as in multi-parameter
PH --- do not admit such a decomposition in general.

Practical invariants for multi-parameter PH include: the rank invariant
\citep{carlsson2009theory}, which is the complete discrete invariant for
2-parameter modules; the Hilbert function, which gives graded Betti numbers
across the bifiltration grid; and the signed barcode/signed measure
\citep{botnan2022signed, loiseaux2023framework}, which decomposes the
persistence module into signed interval summands.

The \texttt{multipers} library \citep{loiseaux2024multipers} implements these
invariants with a scikit-learn-style API and \texttt{gudhi} backend. We use
signed measures as the primary invariant and compute Hilbert functions and
one-dimensional slices of the rank invariant as complementary diagnostics
(Appendix~\ref{app:hilbert-rank}).

\subsection{Why Signed Measures?}

The signed measure $\mu_k$ for homology degree $k$ is a discrete signed
measure on the bifiltration parameter space that encodes where topological
features are born (positive weight) and die (negative weight). For applied
analysis, it has two advantages over the Hilbert function and rank invariant:
(1)~it is computationally the least expensive of the three, and (2)~it
admits a natural income-stratification decomposition by summing absolute
weights within income bands.

% ===================================================================
\section{Data and Methods}
\label{sec:methods}

\subsection{Data: BHPS/USoc Panel Trajectories}

\textbf{Source:} British Household Panel Survey (1991--2008) and
Understanding Society (2009--2023). \textbf{Sample:} $N = 27{,}280$
individuals with $\geq 5$ waves of employment state observations, from the
Paper~1 pipeline. \textbf{Embedding:} 4-gram frequency embedding followed by
PCA to 20 dimensions; each individual is a point in $\mathbb{R}^{20}$.

\textbf{Income variable:} Mean income score derived from employment state
sequences. Each state's final character encodes income band (L, M, H),
mapped to a continuous score: $L = 0, M = 1, H = 2$, averaged across
observed waves. The resulting per-person scalar $y_i \in [0, 2]$ serves as
the second filtration parameter.

\textit{Note:} This proxy income score tracks the trajectory-level income
component rather than a survey-derived monetary income variable. For the
methodological contribution of this paper --- demonstrating that bifiltration
PH reveals income-stratified topology --- the proxy is adequate. A future
analysis (Paper~5) will use equivalised net household income from the USoc
derived variable \texttt{fihhmnnet1\_dv}.

\subsection{Bifiltration Construction}

We construct a two-parameter filtration $\{K_{\varepsilon, \tau}\}$ where:

\textbf{Parameter~1 --- Rips distance ($\varepsilon$):} Standard
Vietoris--Rips filtration on the PCA embedding. An edge between individuals
$i$ and $j$ appears when $\|x_i - x_j\| \leq \varepsilon$.

\textbf{Parameter~2 --- Income threshold ($\tau$):} A vertex $i$ enters the
complex when its mean income $y_i \leq \tau$. This builds the complex upward
from the poorest trajectories.

Formally:
\begin{equation}
  K_{\varepsilon, \tau}
  = \VR_\varepsilon\!\bigl(\{x_i \mid y_i \leq \tau\}\bigr).
  \label{eq:bifiltration}
\end{equation}

\textbf{Landmarking:} We subsample 2,000 landmarks via greedy maxmin
subsampling (seed = 42) from the 27,280-point embedding, preserving
geometric coverage. The Rips radius is set at the 10th percentile of
pairwise distances in a 500-point subsample
($\varepsilon_\text{max} = 11.126$), yielding 215,958 simplices.

\textbf{Income discretisation:} The income axis $\tau$ is discretised using a
\textbf{quantile grid}: 20 thresholds at uniform quantiles of the
2,000-landmark income distribution. Exploratory comparison
(Appendix~\ref{app:discretisation}) confirmed that quantile grids produce
more stable $\HH{1}$ low-income fraction estimates than fixed
(uniform-spacing) grids: convergence to ${\sim}0.33$ by 10 grid steps (range
$[0.33, 0.44]$) versus $[0.26, 0.51]$ for fixed grids. The Rips axis is
discretised into 20 uniform steps from 0 to $\varepsilon_\text{max}$.

\subsection{Multi-Parameter Persistence Computation}

The bifiltration is constructed via
\texttt{multipers.filtrations.RipsLowerstar}, which builds a 2-parameter
simplex tree from the point cloud and income function. The \textbf{signed
measure}~$\mu_k$ for homology degree $k \in \{0, 1\}$ is computed on the
$20 \times 20$ quantile grid using \texttt{multipers.signed\_measure()}.

We report three summary statistics:

\begin{enumerate}
  \item \textbf{Income-quartile mass decomposition:} The $\tau$ axis is
    partitioned into four quartiles (using income percentiles of the full
    landmark set). Total absolute signed measure mass
    $\sum_{(e,t) \in Q_j} |\mu_k(e, t)|$ is computed for each quartile
    $j \in \{1, 2, 3, 4\}$.
  \item \textbf{$\HH{1}$ low-income fraction:}
    $f_\text{low} = \sum_{\tau \leq \text{med}} |\mu_1|\, /\, \sum |\mu_1|$.
  \item \textbf{Q1/Q3 mass ratio:} $r = m_{Q1} / m_{Q3}$, comparing
    topological complexity in the poorest and second-richest quartiles. Under
    income-exchangeability, this should be close to~1.0.
\end{enumerate}

\subsection{Complementary Invariants}
\label{sec:methods-complementary}

We compute two additional invariants using
\texttt{multipers.signed\_measure()} with the \texttt{invariant} parameter:

\textbf{Hilbert function} (\texttt{invariant="hilbert"}): The Hilbert
function $\mathcal{H}_k(\varepsilon, \tau) = \dim H_k(K_{\varepsilon,\tau})$
counts features \emph{alive} at each grid point, providing a cumulative
``state'' view complementary to the signed measure's ``differential''
(birth/death) view. We compute $\mathcal{H}_0$ and $\mathcal{H}_1$ on both
$10 \times 10$ and $20 \times 20$ grids.

\textbf{Rank invariant} (\texttt{invariant="rank"}): The rank invariant
captures which features persist across parameter ranges, returning
4-dimensional support points $(b_\varepsilon, b_\tau, d_\varepsilon,
d_\tau)$ recording the birth and death coordinates of each feature in the
bifiltration. We compute the $\HH{1}$ rank invariant to test whether
features born at low income survive to high income or are confined to
low-income regions.

\subsection{Permutation Null Testing}

\textbf{Null hypothesis:} The income-stratified distribution of $\HH{1}$
signed measure mass is fully explained by the distance structure --- shuffling
income labels does not change the income-quartile decomposition.

\textbf{Procedure:} For $b = 1, \ldots, 999$:
\begin{enumerate}
  \item Shuffle income labels across the 2,000 landmarks (breaking the
    income-trajectory association while preserving the distance structure).
  \item Reconstruct the bifiltration with shuffled incomes.
  \item Compute the $\HH{1}$ signed measure on the \emph{same} quantile grid
    (derived from the original, unshuffled incomes).
  \item Record $f_\text{low}^{(b)}$ and $r^{(b)}$.
\end{enumerate}

\textbf{\pval-value:} Two-sided:
\begin{equation}
  \hat{p} = B^{-1} \sum_{b=1}^{B}
  \mathbb{1}\!\left[
    \bigl|T^{(b)} - \bar{T}\bigr|
    \geq \bigl|T_\text{obs} - \bar{T}\bigr|
  \right],
  \label{eq:pvalue}
\end{equation}
where $\bar{T}$ is the null mean and $T$ is each test statistic.

The permutation preserves the distance matrix (same points in
$\mathbb{R}^{20}$) and the income marginal distribution (same set of income
values), but destroys the association between trajectory position and income
level. A significant result means that the observed income-topology
association is not an artefact of the embedding geometry alone.

\subsection{Comparison to Paper~1 Single-Parameter Baseline}

We compute standard VR PH (ignoring income) on three subsets of the same
2,000 landmarks: the full set, the low-income subset ($y_i \leq Q_1$), and
the high-income subset ($y_i \geq Q_3$). If single-parameter PH shows
similar $\HH{1}$ density across income subsets, this confirms that the income
stratification is only visible through the bifiltration.

% ===================================================================
\section{Results}
\label{sec:results}

\subsection{Bifiltration Structure}

The bifiltration on 2,000 maxmin landmarks from the $27{,}280 \times 20$D
Paper~1 embedding yielded 215,958 simplices at Rips radius 11.126. Landmark
income scores had mean 1.02, median 1.00, and standard deviation 0.48,
centred around the midpoint of the $[0, 2]$ range.

The $\HH{0}$ signed measure (229 support points on the quantile grid) showed
roughly balanced income distribution (low-income fraction = 0.512). Connected
components form at all income levels --- the large-scale cluster structure of
trajectory space is income-symmetric.

The $\HH{1}$ signed measure (218 support points) showed pronounced income
asymmetry: low-income fraction = 0.329, far below the 0.50 expected under
income-exchangeability and below the $\HH{0}$ baseline of 0.51.

\subsection{Income-Quartile Decomposition of $\HH{1}$ Topology}

The $\HH{1}$ signed measure mass decomposes across income quartiles as shown
in Table~\ref{tab:quartile}.

\begin{table}[h]
\centering
\caption{$\HH{1}$ signed measure mass by income quartile ($20 \times 20$
  quantile grid, 2,000 landmarks).}
\label{tab:quartile}
\begin{tabular}{llrr}
  \toprule
  Quartile & Range (income score) & $|\HH{1}$~mass$|$ & Share \\
  \midrule
  Q1 (poorest) & $0.00$--$0.60$ & 20,336  &  9.6\% \\
  Q2           & $0.60$--$1.00$ & 35,448  & 16.7\% \\
  Q3           & $1.00$--$1.36$ & 78,674  & 37.1\% \\
  Q4 (richest) & $1.36$--$2.00$ & 77,524  & 36.6\% \\
  \bottomrule
\end{tabular}
\end{table}

Quartiles Q3 and Q4 carry 73.7\% of total $\HH{1}$ mass. The Q1/Q3 ratio is
0.258 --- topological cycles are approximately four times more prevalent at
Q3 income levels than at Q1.

\textbf{Interpretation:} Mid-to-high income trajectory space has rich
topological structure: multiple career pathways (e.g., different combinations
of employment, self-employment, reduced hours) coexist and form loops in the
embedding. Low-income trajectory space is topologically simpler --- fewer
distinct pathways create less cyclical connectivity. This topological
simplicity reflects the constrained option set that characterises
structurally trapped trajectories.

\subsection{Permutation Null Test Results}

The 999 income-label permutations took 24 minutes on an 8-core i7
(parallelised via \texttt{joblib}). Results are shown in
Table~\ref{tab:null}.

\begin{table}[h]
\centering
\caption{Permutation null test results (999 income-label shuffles).}
\label{tab:null}
\begin{tabular}{lrrr}
  \toprule
  Statistic & Observed & Null mean & \pval-value \\
  \midrule
  $\HH{1}$ low-income fraction & 0.329 & 0.265 & $< 0.001$ \\
  Q1/Q3 mass ratio              & 0.258 & 0.161 & $< 0.001$ \\
  \bottomrule
\end{tabular}
\end{table}

Both test statistics are significantly different from the permutation null
($p < 0.001$). The direction of the effect is notable: the observed
low-income fraction (0.329) is \emph{higher} than the null mean (0.265), and
the observed Q1/Q3 ratio (0.258) is higher than the null mean (0.161).

\textbf{Interpretation:} Under random income assignment, $\HH{1}$ mass
concentrates even more extremely at high $\tau$ than it does with real income
labels. The real income-trajectory association \emph{moderates} the asymmetry
--- low-income trajectories contribute more $\HH{1}$ complexity than random
assignment would predict. This is consistent with the existence of genuine
topological structure (cycling patterns, closed loops) at low income that is
partially masked by the dominant high-income signal.

\begin{figure}[t]
\centering
\includegraphics[width=\textwidth]{null_distribution.png}
\caption{Permutation null distributions for $\HH{1}$ low-income fraction
  (left) and Q1/Q3 mass ratio (right). Red dashed lines mark the observed
  values; grey dotted lines mark the expected values under
  income-exchangeability. The observed $\HH{1}$ low-income fraction (0.329)
  falls in the upper tail of the null distribution ($p < 0.001$).}
\label{fig:null}
\end{figure}

\subsection{Paper~1 Single-Parameter Baseline}

Standard VR PH (single filtration parameter, ignoring income) on the same
2,000 landmarks produced the results in Table~\ref{tab:baseline}.

\begin{table}[h]
\centering
\caption{Single-parameter VR PH on income-stratified subsets of the same
  2,000 landmarks.}
\label{tab:baseline}
\begin{tabular}{lrrr}
  \toprule
  Subset & $N$ & $\HH{1}$ count & $\HH{1}$ density \\
  \midrule
  Full               & 2,000 & 2,587 & 1.294 \\
  Low ($\leq Q_1$)   &   508 &   481 & 0.947 \\
  High ($\geq Q_3$)  &   521 &   460 & 0.883 \\
  \bottomrule
\end{tabular}
\end{table}

$\HH{1}$ density is nearly identical across income subsets (0.95 vs 0.88 per
point). \textbf{Single-parameter PH cannot detect the income stratification}
that the bifiltration reveals. The 73.7\% mass concentration at Q3/Q4 in the
bifiltration is invisible to subset-level VR computation because
single-parameter PH changes the complex when subsetting --- it cannot
decompose the joint topology.

\begin{figure}[t]
\centering
\includegraphics[width=\textwidth]{baseline_comparison.png}
\caption{Paper~1 baseline comparison. Left: $\HH{1}$ feature counts by
  income subset. Right: $\HH{1}$ density (features per point), showing
  negligible income gradient under single-parameter PH.}
\label{fig:baseline}
\end{figure}

\subsection{GMM Regime Correspondence}

The 7-component GMM from Paper~1 was applied to the 2,000 landmarks.
Three regimes of note are shown in Table~\ref{tab:regime}.

\begin{table}[h]
\centering
\caption{Selected GMM regimes and income statistics on the landmark set.}
\label{tab:regime}
\begin{tabular}{lrrr}
  \toprule
  Regime & $N$ landmarks & Mean income & \% below median \\
  \midrule
  R0 (dominant)    & 1,720 & 1.05 & 50.4\% \\
  R5 (high-income) &    16 & 1.55 &  6.2\% \\
  R6 (low-income)  &    99 & 0.47 & 88.9\% \\
  \bottomrule
\end{tabular}
\end{table}

Regime R6 --- with 88.9\% of members below the median income and a mean
income score of 0.47 --- corresponds to the poverty trap cluster. The GMM
regime decomposition is consistent with the $\HH{1}$ mass gradient:
trajectories in R0 and R5 (moderate-to-high income, topologically rich
regions) dominate $\HH{1}$ mass, while R6 trajectories occupy the
topologically simpler low-income region.

\begin{figure}[t]
\centering
\includegraphics[width=\textwidth]{regime_income.png}
\caption{Paper~1 GMM regime income distributions. Left: mean income score by
  regime. Right: proportion of regime members below median income. Regime R6
  is the poverty trap cluster (88.9\% below median).}
\label{fig:regime}
\end{figure}

\subsection{Hilbert Function and Rank Invariant}
\label{sec:results-hilbert-rank}

Complementary invariants computed on $10 \times 10$ and $20 \times 20$ grids
(detailed in Appendix~\ref{app:hilbert-rank}) confirm and extend the signed
measure findings.

\textbf{Hilbert function.} The $\HH{0}$ Hilbert function is
income-symmetric (low-income fraction = 0.512 at both resolutions),
confirming that cluster structure is balanced across the income distribution.
The $\HH{1}$ Hilbert function mirrors the signed measure asymmetry with the
same low-income fraction (0.329), showing a clear gradient from near-zero
mass below the income median to concentrated mass in the upper-right corner
of the $(\varepsilon, \tau)$ plane (Figure~\ref{fig:hilbert}).

\begin{figure}[t]
\centering
\includegraphics[width=\textwidth]{appendix_b/hilbert_heatmap_20x20.png}
\caption{Hilbert function heatmaps at $20 \times 20$ resolution.
  Left: $\HH{0}$ --- income-symmetric (low-income fraction = 0.512).
  Right: $\HH{1}$ --- mass concentrated above the median income line
  (dashed), with near-zero weight below. The upper-right concentration
  confirms that loops are alive primarily at high Rips scales among
  high-income trajectories.}
\label{fig:hilbert}
\end{figure}

\textbf{Rank invariant.} The $\HH{1}$ rank invariant on the $20 \times 20$
grid yields 218 features, of which \textbf{100\% persist to infinity} --- no
$\HH{1}$ feature born at any income threshold dies within the bifiltration
range (Table~\ref{tab:rank}).

\begin{table}[h]
\centering
\caption{$\HH{1}$ rank invariant income persistence ($20 \times 20$ grid).}
\label{tab:rank}
\begin{tabular}{lrr}
  \toprule
  Statistic & Count & Weight fraction \\
  \midrule
  Total features                           & 218 & 1.000 \\
  Born at low income ($\tau \leq$ median)  & 108 & 0.329 \\
  Survive to high income (of born-low)     & 108 & 1.000 \\
  Confined to low income                   &   0 & 0.000 \\
  Born at Q1                               &  53 & ---   \\
  Q1-born surviving to Q3                  &  53 & 1.000 \\
  Persist to infinity                      & 218 & 1.000 \\
  \bottomrule
\end{tabular}
\end{table}

This is a striking negative result. There are no income-confined topological
features in the $\HH{1}$ bifiltration. Every loop that enters the complex at
some income threshold $\tau_0$ remains alive as $\tau$ increases to the
maximum. The income gradient in $\HH{1}$ mass is entirely driven by
\emph{where features are born} --- predominantly at high income levels ---
rather than by features dying at specific income thresholds.

\begin{figure}[t]
\centering
\includegraphics[width=\textwidth]{appendix_b/rank_invariant_20x20.png}
\caption{$\HH{1}$ rank invariant analysis ($20 \times 20$ grid). Left:
  birth vs death income scatter --- empty because all features persist to
  infinity ($d_\tau = \infty$). Right: income persistence summary
  showing that 100\% of features born at low income survive to high income,
  and 100\% of all features persist indefinitely.}
\label{fig:rank}
\end{figure}

\textbf{Interpretation for poverty traps:} The absence of income-confined
features means that the income gradient is monotonic rather than
threshold-based. There is no critical income level at which loops appear or
disappear. Instead, the topology enriches monotonically as the income
threshold increases, with each new income level adding features that persist
indefinitely. Low-income trajectory space is not characterised by distinct
topological structures that could be ``broken through''; rather, it is
characterised by the \emph{absence} of the diverse pathway structures that
emerge at higher income levels.

% ===================================================================
\section{Discussion}
\label{sec:discussion}

\subsection{Topological Simplicity as a Poverty Trap Signature}

The finding inverts a naive expectation. Rather than poverty traps
manifesting as topological \emph{features} --- persistent loops or voids at
low income --- they manifest as topological \emph{absence}. Low-income
trajectory space is structurally simpler: fewer $\HH{1}$ cycles, less
multi-pathway connectivity. This makes formal topological sense: a poverty
trap constrains individuals to a narrow set of employment pathways (cycling
between unemployment, inactivity, and low-pay work), producing a
low-dimensional, acyclic region of the embedding.

High-income trajectory space, by contrast, is topologically diverse: multiple
career structures (full-time $\times$ self-employed, part-time $\times$
employed, career changes across sectors) create higher-dimensional loops in
the embedding. These loops represent the career flexibility that poverty
traps deny.

The rank invariant sharpens this interpretation. Because all $\HH{1}$
features persist to infinity, the topology enriches monotonically with income
--- there is no ``threshold'' beyond which the structure qualitatively
changes. This is more consistent with a continuous gradient of career
opportunity than with a sharp poverty trap boundary. The structural
constraint at low income is not a barrier that could be crossed, but rather
the absence of pathway diversity that accumulates incrementally at higher
income levels.

The finding is consistent with \citeauthor{barrett2013economics}'s
(\citeyear{barrett2013economics}) distinction between structural and
stochastic traps. Structural traps --- where the constraint is on the set of
reachable states --- would produce exactly the topological simplicity we
observe: a restricted pathway set with fewer cycles.

\subsection{The Permutation Result: Real Structure Beyond Geometry}

The permutation null test confirms that the income-topology association is
not a geometric artefact. Under random income assignment, $\HH{1}$ mass
concentrates even more at high $\tau$ (null low-income fraction = 0.265 vs
observed = 0.329). The real income labels \emph{increase} low-income
$\HH{1}$ mass relative to the null by approximately 24\%. This means there
is genuine topological structure at low income --- cycling patterns among
poverty-trap trajectories --- that partially offsets the dominant high-income
signal.

This is a two-part finding:
\begin{enumerate}
  \item \textbf{Global asymmetry:} High-income trajectory space is
    inherently topologically richer (driven by the distance structure of the
    embedding).
  \item \textbf{Income-specific moderation:} The association between income
    and trajectory position adds $\HH{1}$ mass at low income, consistent
    with closed cycling patterns among trapped trajectories.
\end{enumerate}

\subsection{What Multi-Parameter PH Adds Beyond Single-Parameter}

The baseline comparison makes the case directly: single-parameter VR PH on
income-stratified subsets shows $\HH{1}$ densities of 0.95 (low) and 0.88
(high) per point --- effectively indistinguishable. The bifiltration's signed
measure reveals a 4:1 Q3:Q1 mass ratio invisible to subset analysis. This is
because single-parameter subsetting changes the simplicial complex (different
points, different edges), while the bifiltration decomposes the \emph{same}
complex across income thresholds.

The Hilbert function and rank invariant add information beyond what the
signed measure provides. The Hilbert function gives a ``state'' view --- how
many features are alive at each point --- that visually confirms the signed
measure's ``differential'' pattern. The rank invariant answers a question the
signed measure cannot: do features born at one income level persist to
another? The answer --- universal persistence to infinity --- constrains the
class of theoretical models that could generate the observed topology.

\subsection{Methodological Contribution}

This paper demonstrates that multi-parameter persistent homology, via the
\texttt{multipers} library's signed measure, Hilbert function, and rank
invariant, is a tractable analytical tool for panel survey data at the scale
of major UK longitudinal studies. The pipeline --- landmarking, bifiltration
construction, signed measure computation, income-quartile decomposition,
permutation null testing, Hilbert/rank diagnostics --- runs in under 30
minutes on commodity hardware (i7/32\,GB, no GPU). The quantile grid
discretisation produces stable results across resolution settings.

\subsection{Limitations}

\begin{enumerate}
  \item \textbf{Income proxy rather than monetary income:} The L/M/H income
    score derived from employment states is a trajectory-level proxy.
    Analysis with equivalised net household income (available in USoc) is
    planned for Paper~5.

  \item \textbf{Landmark subsampling:} 2,000 of 27,280 individuals may miss
    rare trajectory types. Scale sensitivity analysis at 5,000 landmarks is
    planned before submission.

  \item \textbf{Income discretisation:} The continuous $\tau$ parameter is
    approximated on a 20-step quantile grid. Exploratory analysis confirms
    stability across 10--30 steps, but effects at finer resolution cannot be
    ruled out.

  \item \textbf{Two parameters only:} The bifiltration uses distance $\times$
    income. Geography, education, or cohort as additional filtration
    parameters could reveal further structure but require higher-parameter
    methods with steeper computational cost.

  \item \textbf{Single cohort pool:} The pooled 1991--2023 sample may
    conflate secular trends. Cohort-stratified bifiltration is a natural
    extension.

  \item \textbf{Rank invariant resolution:} The rank invariant was computed
    on discrete grids ($10 \times 10$ and $20 \times 20$). Features that are
    born and die between grid points would appear to persist to infinity.
    Finer-resolution grids or exact (non-gridded) computation may reveal
    income-confined features below the current resolution.
\end{enumerate}

\subsection{Future Directions}

\begin{itemize}
  \item \textbf{Paper~5:} Apply the bifiltration pipeline to SOEP (Germany),
    PSID (US), and CNEF harmonised data to compare income-stratified
    trajectory topology across welfare-state regimes.
  \item \textbf{Scale sensitivity:} 5,000-landmark run before submission.
  \item \textbf{Real income:} Comparison with equivalised net household
    income from USoc (Paper~5 scope).
  \item \textbf{Policy relevance:} If low-income topological simplicity
    reflects constrained career pathways, interventions that \emph{diversify}
    available pathways (training, relocation support, flexible employment
    schemes) would increase topological complexity at low income --- a
    measurable topological outcome.
\end{itemize}

% ===================================================================
\section{Conclusion}
\label{sec:conclusion}

Multi-parameter persistent homology, applied to a bifiltration indexed by
Rips distance and income threshold, reveals that the topological structure of
UK employment trajectories is significantly income-stratified. $\HH{1}$
topological features --- loops representing cyclical career structures ---
concentrate at mid-to-high income levels, with quartiles Q3 and Q4 carrying
approximately four times the signed measure mass of Q1. Permutation null
testing (999 income-label shuffles, $p < 0.001$) confirms that this is not a
geometric artefact, and that real income labels increase low-income $\HH{1}$
mass relative to the null by 24\%.

The Hilbert function confirms this asymmetry with a complementary cumulative
view, while the rank invariant shows that all $\HH{1}$ features persist to
infinity along the income axis --- ruling out income-confined topological
structures and indicating a monotonic enrichment of topology with rising
income.

We interpret the topological simplicity of low-income trajectory space as a
formalisation of structural poverty traps: limited career diversity
constrains individuals to a narrow set of pathways, producing less
topological complexity. The rank invariant's finding of universal persistence
further specifies this: the poverty trap is not a threshold to be crossed but
an absence of career pathway diversity that accumulates only at higher income
levels.

% ===================================================================
\appendix

\section{Discretisation Strategy Comparison}
\label{app:discretisation}

Quantile vs fixed (uniform-spacing) grids were compared on 2,000 maxmin
landmarks. Both grids produce the same qualitative finding ($\HH{1}$ mass
concentrates at upper income quartiles), but quantile grids show superior
stability:

\begin{itemize}
  \item $\HH{1}$ low-income fraction converges to ${\sim}0.33$ by 10 grid
    steps (range $[0.33, 0.44]$) vs fixed grid (range $[0.26, 0.51]$).
  \item Quantile grid produces balanced $\HH{0}$ stratification (51/49) vs
    fixed grid's high-income bias (45/55).
\end{itemize}

% ===================================================================
\section{Hilbert Function and Rank Invariant Analysis}
\label{app:hilbert-rank}

\subsection{Hilbert Function}

The Hilbert function $\mathcal{H}_k(\varepsilon, \tau)
= \dim H_k(K_{\varepsilon, \tau})$ records the number of $k$-dimensional
features alive at each bifiltration grid point. Unlike the signed measure,
which captures birth/death differentials, the Hilbert function provides a
cumulative ``state'' view.

\textbf{$\HH{0}$ Hilbert function.} Total mass = 4,074 ($20 \times 20$) /
4,020 ($10 \times 10$). Income distribution is symmetric (low-income
fraction = 0.512 at both resolutions). The heatmap shows strong positive
weights at low $\varepsilon$ (many components before merging), turning
negative at intermediate $\varepsilon$ --- the standard Rips signature of a
point cloud merging into a single component. This pattern is uniform across
the income axis. Quartile decomposition at $20 \times 20$:
Q1 = 1,002, Q2 = 840, Q3 = 1,241, Q4 = 991.

\textbf{$\HH{1}$ Hilbert function.} Total mass = 211,982 (stable across
resolutions). Low-income fraction = 0.329 (identical to the signed measure).
The heatmap (Figure~\ref{fig:hilbert}) shows near-zero mass below the income
median and concentrated mass in the upper-right corner (high $\varepsilon$,
high $\tau$), consistent with loops being alive predominantly at large Rips
scales among high-income trajectories. Quartile decomposition:
Q1 = 20,336 (9.6\%),
Q2 = 35,448 (16.7\%),
Q3 = 78,674 (37.1\%),
Q4 = 77,524 (36.6\%).

\textbf{Resolution stability.} Hilbert function results are stable between
$10 \times 10$ and $20 \times 20$ grids: total mass is identical
(211,982), and low-income fraction is 0.329 at both resolutions. The
$20 \times 20$ grid provides finer spatial detail in the heatmap without
changing the income-stratification summary.

\subsection{Rank Invariant}

The rank invariant $\rho(a, b) = \text{rank}(H_k(K_a) \to H_k(K_b))$
captures which features persist from grid point $a$ to grid point~$b$. In
\texttt{multipers}, the rank invariant signed measure returns 4D support
points $(b_\varepsilon, b_\tau, d_\varepsilon, d_\tau)$ recording the birth
and death coordinates in the bifiltration.

\textbf{$\HH{1}$ rank invariant ($20 \times 20$).} 218 features, total
weight = 211,982. All 218 features have death coordinates at infinity
($d_\varepsilon = \infty$, $d_\tau = \infty$). See
Table~\ref{tab:rank} and Figure~\ref{fig:rank} for the full income
persistence breakdown.

No $\HH{1}$ feature is confined to a specific income range. Every cycle that
enters the bifiltration --- regardless of the income threshold at which it
appears --- persists across the full income range. The income gradient is
driven entirely by feature \emph{birth rates} at different income levels, not
by differential feature survival. This indicates monotonic topological
enrichment: each successive income threshold adds new features without
destroying existing ones.

\textbf{$10 \times 10$ comparison.} At $10 \times 10$ resolution, 59
features were resolved (all persisting to infinity, same qualitative
pattern), confirming that the universal persistence finding is not sensitive
to grid resolution.

\textbf{Caveat.} On a discrete grid, features that are born and die between
grid points would appear to persist to infinity. Finer-resolution computation
(e.g., $50 \times 50$) or exact computation on the full simplex tree could in
principle reveal short-lived income-confined features below the current
resolution. We consider this unlikely to change the qualitative finding, given
the stability between $10 \times 10$ and $20 \times 20$, but note it as a
limitation.

% ===================================================================
\section{Computational Notes}
\label{app:computational}

\textbf{Code:}
\texttt{trajectory\_tda/topology/multipers\_bifiltration.py} (module),
\texttt{trajectory\_tda/scripts/p04\_run\_nulls.py} (permutation runner),
\texttt{trajectory\_tda/scripts/p04\_appendix\_b.py} (Hilbert/rank runner).

\textbf{Runtime:} 24 minutes for 999 permutations; $< 1$ minute for
Hilbert/rank at each resolution. All on 8-core i7/32\,GB (no GPU). The
\texttt{multipers} library uses sklearn fallback for distance computation
(KeOps unavailable on Windows). Full results in
\texttt{results/p04\_multipers/}.

\textbf{Completed phases:}
\begin{enumerate}
  \item Library validation --- 13/13 synthetic tests passing
  \item Data pipeline --- income extraction from Paper~1 sequences
  \item Development-scale computation --- 2,000 landmarks, quantile grid
  \item Null testing --- 999 permutations, $p < 0.001$
  \item Paper~1 baseline comparison --- $\HH{1}$ density indistinguishable
    across income subsets
  \item GMM regime overlay --- R6 identified as poverty trap cluster
  \item Hilbert function --- $\HH{0}$ symmetric, $\HH{1}$ asymmetric,
    resolution-stable
  \item Rank invariant --- 100\% persistence to infinity, no
    income-confined features
\end{enumerate}

\textbf{Remaining:}
\begin{itemize}
  \item 5,000-landmark scale sensitivity (before submission)
  \item Real income variable comparison (Paper~5 scope)
\end{itemize}

% ===================================================================
% Bibliography placeholder
\begin{thebibliography}{99}

\bibitem[Azariadis and Stachurski(2005)]{azariadis2005poverty}
Azariadis, C. and Stachurski, J. (2005).
\newblock Poverty traps.
\newblock In \emph{Handbook of Economic Growth}, vol.~1, pp.~295--384.

\bibitem[Barrett and Carter(2013)]{barrett2013economics}
Barrett, C.~B. and Carter, M.~R. (2013).
\newblock The economics of poverty traps and persistent poverty: empirical
  and policy implications.
\newblock \emph{Journal of Development Studies}, 49(7):976--990.

\bibitem[Botnan et~al.(2022)]{botnan2022signed}
Botnan, M.~B., Oppermann, S., and Oudot, S. (2022).
\newblock Signed barcodes for multi-parameter persistence via rank
  decompositions.
\newblock In \emph{Symposium on Computational Geometry (SoCG)}.

\bibitem[Cappellari and Jenkins(2004)]{cappellari2004modelling}
Cappellari, L. and Jenkins, S.~P. (2004).
\newblock Modelling low income transitions.
\newblock \emph{Journal of the Royal Statistical Society: Series~C},
  53(4):611--625.

\bibitem[Carlsson(2009)]{carlsson2009topology}
Carlsson, G. (2009).
\newblock Topology and data.
\newblock \emph{Bulletin of the American Mathematical Society},
  46(2):255--308.

\bibitem[Carlsson and Zomorodian(2009)]{carlsson2009theory}
Carlsson, G. and Zomorodian, A. (2009).
\newblock The theory of multidimensional persistence.
\newblock \emph{Discrete \& Computational Geometry}, 42(1):71--93.

\bibitem[Lesnick and Wright(2022)]{lesnick2022computing}
Lesnick, M. and Wright, M. (2022).
\newblock Computing minimal presentations and {B}etti numbers of 2-parameter
  persistent homology.
\newblock \emph{SIAM Journal on Applied Algebra and Geometry},
  6(2):267--298.

\bibitem[Loiseaux et~al.(2023)]{loiseaux2023framework}
Loiseaux, D., Carri\`{e}re, M., and Blumberg, A.~J. (2023).
\newblock A framework for fast and stable representations of multiparameter
  persistent homology decompositions.
\newblock \emph{Advances in Neural Information Processing Systems (NeurIPS)},
  36.

\bibitem[Loiseaux et~al.(2024)]{loiseaux2024multipers}
Loiseaux, D. et~al. (2024).
\newblock Multipers: multiparameter persistence for machine learning.
\newblock \emph{Journal of Open Source Software}, 9(95):6773.

\bibitem[Shildrick et~al.(2012)]{shildrick2012poverty}
Shildrick, T., MacDonald, R., Webster, C., and Garthwaite, K. (2012).
\newblock \emph{Poverty and Insecurity: Life in Low-Pay, No-Pay Britain}.
\newblock Policy Press.

\bibitem[Vipond(2020)]{vipond2020multiparameter}
Vipond, O. (2020).
\newblock Multiparameter persistence landscapes.
\newblock \emph{Journal of Machine Learning Research}, 21:1--38.

\end{thebibliography}
